
\scanpage{149}
$\mathscr H$ Hilbert. Si $u$ est à trace, $\frac1{1+u}=1-u\frac{R(u)}{\det(1+u)}$

\[
R(u)=\sum_0^\infty P_n(u)=\frac{\det(1+u)}{1+u}
\]

\uline{Théorème} \uline{1) Soit $u$ opérateur \deleted{\unsure{herm.}} à trace normal. Alors}

\[
\boxed{|P_n(u)|\leq P_n(|u|)\ ;\ |R(u)|\leq R(|u|)}\tag*{\unsure{$(*),\,R)$|$(*)$|$(*),\,({**})$}}
\]

En particulier, si $u$ est $\geq0$, il en est de $=$ de $P_n(u)$ et de $R(u)$

2) Si $u\geq0$, alors la suite des termes $\alpha_{n-k}(u)u^k$ dans l'expression de $P_n(u)$, $P_n(u)=\sum_{k=0}^n\alpha_{n-k}(u)u^k$, est une suite \uline{décroissante} d'hermitiens positifs \deleted{\unsure{quand $k$}}.

Donc on a en particulier, pour $u\geq0$

\[
\deleted{\unsure{\alpha_n(u)}}\quad
\boxed{0\leq P_n(u)\leq\alpha_n(u)\mathbin{.}1\qquad\text{d'où}\quad
\deleted{\unsure{R(u)\leq\alpha_n(u)\,1}}\inserted{0\leq R(u)\leq\det(1+u)1}}\tag*{(**)}
\]

\deleted{\unsure{3)}} \uline{Corollaire} Si $u$ est opérateur \deleted{\unsure{de rang}} à trace normal, on a

\[
\left\{\begin{aligned}
\|R(u)\|_\infty&\leq\|R(|u|)\|_\infty\leq\det(1+|u|)\leq e^{\|u\|_1}\\
\|P_n(u)\|_\infty&\leq\|P_n(|u|)\|_\infty\leq\alpha_n(|u|)\leq\frac{\|u\|_1^n}{n!}
\end{aligned}\right.\tag*{(***)}
\]

\uline{Démonstration} On démontre le

\uline{Lemme} Si \inserted{$u$ est opérateur normal} $u=\sum_1^\infty\lambda_i e_i\otimes e'_i$, alors $P_n(u)=\sum\unsure{p_i(u)|\rho_i(u)}e_i\otimes e_i$, où $p_i(u)=\sum_{(I)}'\lambda_{i_1}\cdots\lambda_{i_n}$, $\sum'$ étant étendue à tous les \deleted{\unsure{systèmes}} \deleted{ensembles}\inserted{$I=\{i_1\ldots i_n\}$} de $n$ indices, \deleted{\unsure{qui tels que $I$} \illegible{2}}\inserted{ne contenant pas $i$}

Pour le voir, on calcule successivement (dans la base $(e_i)$) \deleted{\unsure{$u^n$}}$u^n-$$\alpha_1(u)u^{n-1}$, $u^n-\alpha_1(u)u^{n-1}+\alpha_2(u)u^{n-1}$ par récurrence on obtient $\sum_{k=0}^{\unsure{p|n}}u^{n-k}\alpha_k(u)(-1)^k=(-1)^n\sum'\lambda_{i_1}\cdots\lambda_{i_k}$ d'où \deleted{\unsure{$P_n=$ pour tout}} le résultat.

Maintenant 1) devient évident. 2) Se vérifie sur les valeurs propres de façon immédiate.

\uline{Remarque} Notons que si $\rho_i$ sont les valeurs propres de $R(u)$, on a $\rho_i=\prod_{j\neq i}(1+\lambda_j)$, d'où $|\rho_i|\leq\prod_j(1+|\lambda_j|)=\det(1+\lvert u\rvert)$, on retrouve une
% archive sheet number: 150
